% Generated draft: requires mathematical and source review.
\documentclass[11pt]{article}
\usepackage[margin=25mm]{geometry}
\usepackage{fontspec}
\setmainfont{Latin Modern Roman}
\usepackage{amsmath,amssymb,mathtools,unicode-math}
\setmathfont{Latin Modern Math}
\usepackage{enumitem,needspace,xurl,hyperref}
\setlength{\parindent}{0pt}
\setlength{\parskip}{5pt}
\setlength{\emergencystretch}{3em}
\newcommand{\N}{\mathbb{N}}
\newcommand{\Z}{\mathbb{Z}}
\newcommand{\Q}{\mathbb{Q}}
\newcommand{\R}{\mathbb{R}}
\newcommand{\C}{\mathbb{C}}
\newcommand{\F}{\mathbb{F}}
\newcommand{\A}{\mathbb{A}}
\newcommand{\PP}{\mathbb{P}}
\newcommand{\E}{\mathbb{E}}
\newcommand{\eps}{\varepsilon}
\newcommand{\dd}{\,\mathrm{d}}
\begin{document}
\section*{UnsolvedMath ID: 30003965}
\textbf{Source problem code:} catalogue entry
\subsection{Definitions and mathematical statement}
In the standard interpretation in which \(\Gamma\) is a compact metric graph, a divisor is a finite integral sum \(D=\sum_{i=1}^{n}a_i p_i\), with degree \(\deg D=\sum_{i=1}^{n}a_i=d\). A divisor is effective if all its coefficients are nonnegative. Linear equivalence is generated by divisors of continuous piecewise-linear functions with integral slopes. The divisor rank is \(r_\Gamma(D)=-1\) if no effective divisor is linearly equivalent to \(D\); otherwise,
\[
r_\Gamma(D)=\max\bigl\{k\in\Z_{\geq0}: D-E\text{ is linearly equivalent to an effective divisor for every effective divisor }E\text{ of degree }k\bigr\}.
\]
These definitions specify the compact, unweighted metric-graph convention only; the supplied record does not establish whether that convention is intended, or whether vertex weights or unbounded edges are allowed.

For a divisor \(\widetilde D\) on a smooth projective curve \(X/K\), write
\[
r_X(\widetilde D)=\dim_K H^0\bigl(X,\mathcal O_X(\widetilde D)\bigr)-1.
\]
A specialization map \(\operatorname{sp}\) is understood to arise from a suitable model or skeleton of \(X\), identified with the prescribed tropical curve. The exact model and identification are part of the unspecified meaning of ``dual tropical curve'' in the source.

\textbf{Supplied lifting question.} Given a tropical curve \(\Gamma\) and a divisor
\[
D=\sum_{i=1}^{n}a_i p_i,\qquad \deg D=d,\qquad r_\Gamma(D)=r,
\]
does there exist a non-Archimedean field \(K\), a smooth projective curve \(X/K\) whose dual tropical curve is \(\Gamma\), and a divisor \(\widetilde D\) on \(X\), with multiplicity profile \((a_1,\ldots,a_n)\), such that
\[
\operatorname{sp}(\widetilde D)=D
\qquad\text{and}\qquad
r_X(\widetilde D)=r?
\]
The intended quantification is over every input pair \((\Gamma,D)\) satisfying the stated conditions, with \(K\), \(X\), and \(\widetilde D\) allowed to depend on that pair. This is a question, not an assertion that every pair admits a lift.

The source does not define ``multiplicity profile.'' If the \(p_i\) are distinct support points and the lifted support consists of distinct \(K\)-rational points \(q_i\), its natural expression would be
\[
\widetilde D=\sum_{i=1}^{n}a_i q_i,\qquad \operatorname{sp}(q_i)=p_i.
\]
This expression is conditional on that interpretation, not an additional supplied hypothesis. The record does not specify positivity of the \(a_i\), rationality of lifted points, allowable field extensions, or treatment of residue degrees for closed points. Consequently, a fully determined lifting problem requires these conventions, together with the tropical-curve and specialization conventions, to be fixed.
\subsection{Short English statement}
The supplied question asks whether a tropical divisor can be lifted to a divisor on a smooth projective curve over some non-Archimedean field while preserving its specialization, multiplicity profile, and rank exactly. The meanings of tropical curve, dual tropical curve, and multiplicity profile require clarification before the question has a single precise interpretation.
\subsection{Sources}
\begin{itemize}
\item Local UnsolvedMath registry: permanent ID \(30003965\), title \emph{Lifting Tropical Divisors with Prescribed Rank}, exact source problem code \texttt{catalogue entry}; collection: \emph{Oberwolfach Reports Open Problems}. The supplied record gives proposed year \(2018\) and names no individual proposer.
\item Supplied source citation: \emph{Flat Surfaces and Algebraic Curves} (2019), Oberwolfach Reports, pp.\ 2583--2650, DOI \texttt{10.4171/owr/2018/42}. Supplied source URL and literature-source link labelled ``OWR report'': \url{https://doi.org/10.4171/owr/2018/42}. The supplied background describes these reports as documenting mathematical workshops and their open-problem discussions.
\item Supplied literature source: \emph{Tropical canonical divisors}, \url{https://arxiv.org/abs/1710.06401}.
\item Supplied literature source: \emph{Tropical pluri-canonical divisors}, \url{https://londmathsoc.onlinelibrary.wiley.com/doi/full/10.1112/jlms.70027}.
\item Literature assessment dated exactly \texttt{2026-08-21}: the supplied curation labels the problem open, describes the general realizability problem as difficult by ``Murphy law,'' states that canonical and pluri-canonical cases have criteria, and states that no general solution is known. This is supplied dated curation, not an independently verified literature assessment; the supplied background explicitly requests independent verification of time-sensitive status and cited claims.
\item The supplied registry labels the statement \texttt{corrected\_verified} and reports manual semantic, typographic, and LaTeX review. These are source metadata, not independent verification of the mathematical formulation or its status.
\end{itemize}
\end{document}
